\documentclass{beamer}

\mode<presentation>





\usetheme{Frankfurt}%
\usecolortheme{seagull}
\logo{\includegraphics[height=.25in]{clarksonGreen}}

\definecolor{garnet}{RGB}{136,0,0}
%\definecolor{clarksonGreen}{RGB}{0,71,28}
\definecolor{clarksonGreen}{RGB}{0,52,21}
\setbeamercolor{palette primary}{fg=clarksonGreen,bg=white}
\setbeamercolor{palette secondary}{fg=clarksonGreen,bg=white}
\setbeamercolor{palette tertiary}{fg=clarksonGreen,bg=white}
\setbeamercolor{palette quaternary}{bg=clarksonGreen,fg=white}
\setbeamercolor{block title}{fg=black,bg=black!15}
\setbeamercolor{block body}{fg=black,bg=black!10}
\setbeamercolor{titlelike}{bg=clarksonGreen,fg=white} % parent=palette quaternary}

\newcommand{\half}{\mbox{$\frac{1}{2}$}}
\newcommand{\deltat}{\mbox{$\triangle t$}}
\newcommand{\deltax}{\mbox{$\triangle x$}}
\newcommand{\deltay}{\mbox{$\triangle y$}}

\newcommand{\deriv}[2]{\frac{d}{d#2}#1}
\newcommand{\derivTwo}[2]{\frac{d^2}{d#2^2}#1}

\newcommand{\lp}{\left(}
\newcommand{\rp}{\right)}


\newcommand{\arrayTwo}[4]{
  \left[
  \begin{array}{rr}
    #1 & #2 \\
    #3 & #4
  \end{array}
  \right]
}


\newcommand{\vecTwo}[2]{
  \left[
  \begin{array}{r}
    #1 \\  #2
  \end{array}
  \right]
}


\newcommand{\vecThree}[3]{
  \left[
  \begin{array}{r}
    #1 \\  #2 \\ #3
  \end{array}
  \right]
}


\newcommand{\arrayThree}[9]{
  \left[
    \begin{array}{rrr}
      #1 & #2 & #3 \\
      #4 & #5 & #6 \\
      #7 & #8 & #9
    \end{array}
  \right]
}


\newcommand{\startRowFour}{
  \left[
    \begin{array}{rrr|r}
}

\newcommand{\oneRowFour}[4] {
      #1 & #2 & #3 & #4 \\
}

\newcommand{\stopRowOps}{
    \end{array}
  \right]
}



\begin{document}

\title{Ordinary Differential Equations}
\subtitle{Math 232 - Week 10, Day 3}

\author{Kelly Black}
\institute{Clarkson University}
\date{4 November 2011}

\begin{frame}
  \titlepage
\end{frame}

\begin{frame}
  \frametitle{Outline}
  \tableofcontents[pausesection,hideallsubsections]
\end{frame}


\section{Eigen Vectors and Eigen Values}


\begin{frame}
  \frametitle{Eigen Vectors and Eigen Values}

  Suppose that we have a system of differential equations,
  \begin{eqnarray*}
    \frac{d}{dt} \vec{x} & = & A \vec{x}.
  \end{eqnarray*}
  Assume that $A$ is a square matrix.

  What if we can find a vector, $\vec{v}$, that satisfies
  \begin{eqnarray*}
    A \vec{v} & = & \lambda \vec{v},
  \end{eqnarray*}
  where $\lambda$ is a scalar number?

\end{frame}


\begin{frame}
  If so. then....
  \begin{eqnarray*}
    \frac{d}{dt} \vec{v} & = & A \vec{v}, \\
    & = & \lambda \vec{v}.
  \end{eqnarray*}
  
  This implies that
  \begin{eqnarray*}
    \deriv{x_1}{t} & = & \lambda x_1, \\
    \deriv{x_2}{t} & = & \lambda x_2, \\
    \deriv{x_3}{t} & = & \lambda x_3, \\
    \vdots         &   & \vdots \\
    \deriv{x_n}{t} & = & \lambda x_n
  \end{eqnarray*}

  We can solve each one of these equations separately!

\end{frame}


\begin{frame}
  \frametitle{Eigen Vectors and Eigen Values}

  Definition: Given a square matrix, $A$, an eigen vector, $\vec{v}$, and an
  eigen value, $\lambda$, satisfy
  \begin{eqnarray*}
    A \vec{v} & = & \lambda \vec{v}.
  \end{eqnarray*}

  Given $A$ how can we find them?

\end{frame}


\begin{frame}
  \frametitle{How to find}

  \begin{eqnarray*}
    A \vec{v} & = & \lambda \vec{v}, \\
    A \vec{v} - \lambda \vec{v} & = & 0, \\
    A \vec{v} - \lambda I \vec{v} & = & 0, \\
    \lp A - \lambda I \rp \vec{v} & = & 0.
  \end{eqnarray*}

  One solution is $\vec{v}=\vec{0}$. We do not like this solution. We
  must have a situation where there is not a unique solution. If there
  is not a unique solution then this implies that
  \begin{eqnarray*}
    \det\lp A - \lambda I \rp & = & 0.
  \end{eqnarray*}
  Find a value for $\lambda$ that makes this true. (Note that there
  will be infinite solutions for $\vec{v}$.)

\end{frame}


\begin{frame}
  \frametitle{Finding Eigen Vectors and Eigen Values}

  Given a square matrix, $A$, to find the eigen vectors and eigen
  values do the following:
  \begin{enumerate}
  \item Determine the characteristic equation:
    \begin{eqnarray*}
      \det\lp A - \lambda I \rp & = & 0.
    \end{eqnarray*}
  \item Solve the equation for the values of $\lambda$.
  \item Find the eigen vector associated with each value of $\lambda$,
    \begin{eqnarray*}
      \lp A - \lambda_j I \rp \vec{v}_j & = & \vec{0}.
    \end{eqnarray*}
    (Note that there will be infinite solutions for $\vec{v}$.)
  \end{enumerate}

\end{frame}

\section{Examples}

\begin{frame}
  \frametitle{Example}

  \begin{eqnarray*}
    A & = & \arrayTwo{12}{-10}{15}{-13}.
  \end{eqnarray*}

  Find the characteristic equation:
  \begin{eqnarray*}
    \det\lp\arrayTwo{12-\lambda}{-10}{15}{-13-\lambda}\rp & = & \lambda^2+\lambda-6.
  \end{eqnarray*}

  Find the values of $\lambda$:
  \begin{eqnarray*}
    \lambda^2+\lambda-6 & = & (\lambda-2)(\lambda+3), \\
    & = & 0.
  \end{eqnarray*}
  So $\lambda_1 = -3$ and $\lambda_2=2$.

\end{frame}


\begin{frame}
  Find the vectors:

  $\lambda_1 = 2$:
  \begin{eqnarray*}
    \arrayTwo{10}{-10}{15}{-15} \vecTwo{x}{y} & = & \vecTwo{0}{0}, \\
    10x - 10y & = & 0, \\
    15x - 15y & = & 0, \\
    x & = & y.
  \end{eqnarray*}

  The first eigen vector is 
  \begin{eqnarray*}
    \vec{v}_1 & = & \vecTwo{y}{y}, \\
    & = & y \vecTwo{1}{1}.
  \end{eqnarray*}

  Let
  \begin{eqnarray*}
    \vec{v}_1 & = & \vecTwo{1}{1}.
  \end{eqnarray*}

\end{frame}

\begin{frame}
  Find the vectors:

  $\lambda_2 = -3$:
  \begin{eqnarray*}
    \arrayTwo{15}{-10}{15}{-10} \vecTwo{x}{y} & = & \vecTwo{0}{0}, \\
    15x - 10y & = & 0, \\
    15x - 10y & = & 0, \\
    x & = & \frac{2}{3} y.
  \end{eqnarray*}

  The second eigen vector is 
  \begin{eqnarray*}
    \vec{v}_2 & = & \vecTwo{\frac{2}{3}y}{y}, \\
    & = & y \vecTwo{\frac{2}{3}}{1}.
  \end{eqnarray*}

  Let
  \begin{eqnarray*}
    \vec{v}_2 & = & \vecTwo{2}{3}.
  \end{eqnarray*}

\end{frame}






\begin{frame}
  \frametitle{Example}

  \begin{eqnarray*}
    A & = & \arrayTwo{4}{18}{-3}{-11}.
  \end{eqnarray*}

  Find the characteristic equation:
  \begin{eqnarray*}
    \det\lp\arrayTwo{4-\lambda}{18}{-3}{-11-\lambda}\rp & = & \lambda^2+7\lambda+10.
  \end{eqnarray*}

  Find the values of $\lambda$:
  \begin{eqnarray*}
    \lambda^2+7\lambda+10 & = & (\lambda+2)(\lambda+5), \\
    & = & 0.
  \end{eqnarray*}
  So $\lambda_1 = -2$ and $\lambda_2=-5$.


\end{frame}


\begin{frame}
  Find the vectors:

  $\lambda_1 = -2$:
  \begin{eqnarray*}
    \arrayTwo{6}{18}{-3}{-9} \vecTwo{x}{y} & = & \vecTwo{0}{0}, \\
    6x + 18y & = & 0, \\
    -3x - 9y & = & 0, \\
    x & = & -3y.
  \end{eqnarray*}

  The first eigen vector is 
  \begin{eqnarray*}
    \vec{v}_1 & = & \vecTwo{-3y}{y}, \\
    & = & y \vecTwo{-3}{1}.
  \end{eqnarray*}

  Let
  \begin{eqnarray*}
    \vec{v}_1 & = & \vecTwo{-3}{1}.
  \end{eqnarray*}

\end{frame}

\begin{frame}
  Find the vectors:

  $\lambda_2 = -5$:
  \begin{eqnarray*}
    \arrayTwo{9}{18}{-3}{-6} \vecTwo{x}{y} & = & \vecTwo{0}{0}, \\
    9x + 18y & = & 0, \\
    -3x - 6y & = & 0, \\
    x & = & -2 y.
  \end{eqnarray*}

  The second eigen vector is 
  \begin{eqnarray*}
    \vec{v}_2 & = & \vecTwo{-2y}{y}, \\
    & = & y \vecTwo{-2}{1}.
  \end{eqnarray*}

  Let
  \begin{eqnarray*}
    \vec{v}_2 & = & \vecTwo{-2}{1}.
  \end{eqnarray*}

\end{frame}



\begin{frame}
  \frametitle{Theorem - Distinct Eigen Values Makes Life Easier}

  Theorem: If the eigen values of an $n\times n$ matrix include $n$
  distinct values then the eigen vectors are linearly independent.

  \vfill

  \uncover<2->
  {
    So what? When we start solving DEs it will mean we are essentially
    done. 
  }

  \vfill

\end{frame}



\begin{frame}
  \frametitle{Example}

  See example 5 on page 318.
  \begin{eqnarray*}
    A & = & \arrayThree{-2}{1}{1}{1}{-2}{1}{1}{1}{-2}.
  \end{eqnarray*}

  Find the characteristic equation:
  \begin{eqnarray*}
    \det\lp\arrayThree{-2-\lambda}{1}{1}{1}{-2-\lambda}{1}{1}{1}{-2-\lambda}\rp
    & = & \lambda (\lambda+3)^2.
  \end{eqnarray*}

  Find the values of $\lambda$:
  \begin{eqnarray*}
    \lambda (\lambda+3)^2 & = & 0.
  \end{eqnarray*}
  So $\lambda_1 = 0$ and $\lambda_2=-3$.


\end{frame}


\begin{frame}
  Find the vectors:

  $\lambda_1 = 0$:
  \begin{eqnarray*}
    \arrayThree{-2}{1}{1}{1}{-2}{1}{1}{1}{-2}\vecThree{x}{y}{z} & = & \vecThree{0}{0}{0}, \\
  \end{eqnarray*}

  Express this system in an augmented matrix and put it in reduced row
  echelon form:
  \begin{eqnarray*}
    \mathrm{RREF}\startRowFour
    \oneRowFour{-2}{1}{1}{0} 
    \oneRowFour{1}{-2}{1}{0}
    \oneRowFour{1}{1}{-2}{0}
    \stopRowOps
    & = & 
    \startRowFour
    \oneRowFour{1}{0}{-1}{0} 
    \oneRowFour{0}{1}{-1}{0}
    \oneRowFour{0}{0}{0}{0}
    \stopRowOps
  \end{eqnarray*}


\end{frame}

\begin{frame}
  The first eigen vector is 
  \begin{eqnarray*}
    \vec{v}_1 & = & \vecThree{z}{z}{z}, \\
    & = & z \vecThree{1}{1}{1}.
  \end{eqnarray*}

  Let
  \begin{eqnarray*}
    \vec{v}_1 & = & \vecThree{1}{1}{1}.
  \end{eqnarray*}

\end{frame}

\begin{frame}
  Find the vectors:

  $\lambda_2 = -3$:
  \begin{eqnarray*}
    \arrayThree{1}{1}{1}{1}{1}{1}{1}{1}{1}
    \vecThree{x}{y}{z} & = & \vecThree{0}{0}{0}, \\
  \end{eqnarray*}

  Express this system in an augmented matrix and put it in reduced row
  echelon form:
  \begin{eqnarray*}
    \mathrm{RREF}\startRowFour
    \oneRowFour{1}{1}{1}{0} 
    \oneRowFour{1}{1}{1}{0}
    \oneRowFour{1}{1}{1}{0}
    \stopRowOps
    & = & 
    \startRowFour
    \oneRowFour{1}{1}{1}{0} 
    \oneRowFour{0}{0}{0}{0}
    \oneRowFour{0}{0}{0}{0}
    \stopRowOps
  \end{eqnarray*}

\end{frame}

\begin{frame}

  The second eigen vector is 
  \begin{eqnarray*}
    \vec{v}_2 & = & \vecThree{-y-z}{y}{z}, \\
    & = & y \vecThree{-1}{1}{0} + z \vecThree{-1}{0}{1}.
  \end{eqnarray*}

  Let
  \begin{eqnarray*}
    \vec{v}_2 & = & \vecThree{-1}{1}{0}, \\
    \vec{v}_3 & = & \vecThree{-1}{0}{1}.
  \end{eqnarray*}

\end{frame}



\begin{frame}
  \frametitle{Example}

  \begin{eqnarray*}
    A & = & \arrayThree{2}{1}{1}{0}{1}{1}{0}{0}{1}.
  \end{eqnarray*}

  Find the characteristic equation:
  \begin{eqnarray*}
    \det\lp\arrayThree{2-\lambda}{1}{1}{0}{1-\lambda}{1}{0}{0}{1-\lambda}\rp
    & = & (2-\lambda) (1-\lambda)^2.
  \end{eqnarray*}

  Find the values of $\lambda$:
  \begin{eqnarray*}
    (2-\lambda) (1-\lambda)^2 & = & 0.
  \end{eqnarray*}
  So $\lambda_1 = 2$ and $\lambda_2=1$.


\end{frame}


\begin{frame}
  Find the vectors:

  $\lambda_1 = 2$:
  \begin{eqnarray*}
    \arrayThree{0}{1}{1}{0}{-1}{1}{0}{0}{-1}\vecThree{x}{y}{z} & = & \vecThree{0}{0}{0}, \\
  \end{eqnarray*}

  Express this system in an augmented matrix and put it in reduced row
  echelon form:
  \begin{eqnarray*}
    \mathrm{RREF}\startRowFour
    \oneRowFour{0}{1}{1}{0} 
    \oneRowFour{0}{-1}{1}{0}
    \oneRowFour{0}{0}{-1}{0}
    \stopRowOps
    & = & 
    \startRowFour
    \oneRowFour{0}{1}{0}{0} 
    \oneRowFour{0}{0}{1}{0}
    \oneRowFour{0}{0}{0}{0}
    \stopRowOps
  \end{eqnarray*}


\end{frame}

\begin{frame}
  The first eigen vector is 
  \begin{eqnarray*}
    \vec{v}_1 & = & \vecThree{x}{0}{0}, \\
    & = & x \vecThree{1}{0}{0}.
  \end{eqnarray*}

  Let
  \begin{eqnarray*}
    \vec{v}_1 & = & \vecThree{1}{0}{0}.
  \end{eqnarray*}

\end{frame}

\begin{frame}
  Find the vectors:

  $\lambda_2 = 1$:
  \begin{eqnarray*}
    \arrayThree{1}{1}{1}{0}{0}{1}{0}{0}{0}
    \vecThree{x}{y}{z} & = & \vecThree{0}{0}{0}, \\
  \end{eqnarray*}

  Express this system in an augmented matrix and put it in reduced row
  echelon form:
  \begin{eqnarray*}
    \mathrm{RREF}\startRowFour
    \oneRowFour{1}{1}{1}{0} 
    \oneRowFour{0}{0}{1}{0}
    \oneRowFour{0}{0}{0}{0}
    \stopRowOps
    & = & 
    \startRowFour
    \oneRowFour{1}{1}{0}{0} 
    \oneRowFour{0}{0}{1}{0}
    \oneRowFour{0}{0}{0}{0}
    \stopRowOps
  \end{eqnarray*}

\end{frame}

\begin{frame}

  The second eigen vector is 
  \begin{eqnarray*}
    \vec{v}_2 & = & \vecThree{-y}{y}{0}, \\
    & = & y \vecThree{-1}{1}{0}.
  \end{eqnarray*}

  Let
  \begin{eqnarray*}
    \vec{v}_2 & = & \vecThree{-1}{1}{0}.
  \end{eqnarray*}

  There are only two eigen vectors!

\end{frame}






\end{document}

% LocalWords:  Clarkson pausesection hideallsubsections Eigen eigen
